\section*{Hoofdstuk \ref{ch:b3:group-theory-applied} --- Toegepaste groepentheorie}

\begin{solution}{exo:b3:group-theory-applied:1}
$n_{A_1} = \frac14(5 - 1 + 3 + 1) = 2$, $n_{A_2} = \frac14(5 - 1 - 3 - 1) = 0$,
$n_{B_1} = \frac14(5 + 1 + 3 - 1) = 2$, $n_{B_2} = \frac14(5 + 1 - 3 + 1) = 1$:
$\Gamma = 2A_1 + 2B_1 + B_2$, van dimensie 5.
\end{solution}

\begin{solution}{exo:b3:group-theory-applied:2}
Dezelfde telling van atomen op hun plaats als bij water: $\Gamma_{\mathrm{vib}} = 2A_1 + B_1$. $A_1$ ($z$) en
$B_1$ ($x$) zijn \omterm{def:b3:group-theory-applied:activity}{IR-actief}, en beide zijn ook \omterm{def:b3:group-theory-applied:activity}{Raman-actief} ($x^2$, $xz$):
drie banden in elk spectrum.
\end{solution}

\begin{solution}{exo:b3:group-theory-applied:3}
$B_1\otimes B_2$: \omterm{def:b3:point-groups:representation}{karakters} $(1, 1, -1, -1)$ $= A_2$. $A_2\otimes B_1$: $(1, -1, -1,
1) = B_2$. In $C_{3v}$ heeft $E\otimes E$ de \omterm{def:b3:point-groups:representation}{karakters} $(4, 1, 0)$; reductie:
$A_1 + A_2 +
E$.
\end{solution}

\begin{solution}{exo:b3:group-theory-applied:4}
Infrarood: de twee $t_2$-trillingen, 3019 en \qty{1306}{cm^{-1}}. Raman: alle
vier ($a_1$,
$e$, $t_2$ transformeren alle als kwadratische functies). Geen
\omterm{def:b3:point-groups:inversion}{inversiecentrum}, dus samenvallende banden zijn toegestaan.
\end{solution}

\begin{solution}{exo:b3:group-theory-applied:5}
Atomen op hun plaats $(4, 1, 2, 4, 1, 2)$ maal $\chi_{xyz} = (3, 0, -1, 1, -2, 1)$ geeft
$\chi_{3N} = (12, 0, -2, 4, -2, 2)$, wat reduceert tot $A_1' + A_2' + 3E' + 2A_2'' +
E''$. Trek $E' + A_2''$ (translaties) en $A_2' + E''$ (rotaties) af: $A_1' + 2E'
+ A_2''$. IR: $2E' + A_2''$, drie banden; Raman: $A_1' + 2E'$, drie banden. $A_2''$
(buigvibratie uit het vlak) is alleen \omterm{def:b3:group-theory-applied:activity}{IR-actief}, $A_1'$ (symmetrische
strekvibratie) alleen \omterm{def:b3:group-theory-applied:activity}{Raman-actief}, de $E'$-trillingen verschijnen in beide.
\end{solution}

\begin{solution}{exo:b3:group-theory-applied:6}
IR: de twee $T_{1u}$-trillingen, twee banden. Raman: $A_{1g}$, $E_g$, $T_{2g}$,
drie banden.
$T_{2u}$ is stil. \ce{SF6} heeft een \omterm{def:b3:point-groups:inversion}{inversiecentrum}: wederzijdse uitsluiting.
\end{solution}

\begin{solution}{exo:b3:group-theory-applied:7}
Onder $E$, $C_3$, $C_3^2$ gaat de functie $h_2$ naar $h_2$, $h_3$, $h_1$ (\omterm{def:b3:point-groups:representation}{karakters} 2,
$-1$, $-1$); spiegelingen dragen 0 bij. $\hat P_Eh_2 \propto 2h_2 - h_3 - h_1$.
De overlap ervan met $2h_1 - h_2 - h_3$ (atomaire overlappen verwaarloosd)
is $-3$, en deze laatste heeft norm $\sqrt6$; na aftrek van de projectie:
$(2h_2 - h_1 - h_3) +
\frac12(2h_1 - h_2 - h_3) = \frac32(h_2 - h_3)$.
\end{solution}

\begin{solution}{exo:b3:group-theory-applied:8}
fac ($C_{3v}$, \omterm{def:b3:point-groups:representation}{karakters} $3, 0, 1$): $A_1 + E$, twee IR-banden. mer ($C_{2v}$,
\omterm{def:b3:point-groups:representation}{karakters} $3, 1, 3, 1$): $2A_1 + B_1$, drie banden. \ce{M(CO)6} ($O_h$): $A_{1g} + E_g
+ T_{1u}$, één IR-band ($T_{1u}$).
\end{solution}

\begin{solution}{exo:b3:group-theory-applied:9}
\omterm{def:b3:point-groups:representation}{Karakters} $6, 0, 0, 2, 2, 0, 0, 0, 4, 2$: alleen $E$, de $C_4$ en $C_2$ langs
de assen (twee liganden op hun plaats), de drie $\sigma_h$ (vier) en de zes
$\sigma_d$ (twee) laten liganden op hun plaats. De reductie geeft $A_{1g} + E_g + T_{1u}$. De
metaalorbitalen $s$
($a_{1g}$), $d_{z^2}$ en $d_{x^2-y^2}$ ($e_g$), en $p_x$, $p_y$, $p_z$ ($t_{1u}$) combineren;
$t_{2g}$ blijft niet-bindend.
\end{solution}

\begin{solution}{exo:b3:group-theory-applied:10}
$n_i = \frac1h(1\cdot h\cdot\chi_i(E)) = d_i$. De dimensie van de reguliere
representatie is $h$, en ook $\sum_in_id_i = \sum_id_i^2$.
\end{solution}

\begin{solution}{exo:b3:group-theory-applied:11}
In $D_{2h}$ (molecuul langs $z$): atomen op hun plaats $(4, 4, 0, 0, 0, 0, 4, 4)$, $\chi_{3N}
= (12, -4, 0, 0, 0, 0, 4, 4)$, wat reduceert tot $2A_g + 2B_{2g} + 2B_{3g} + 2B_{1u} + 2B_{2u}
+ 2B_{3u}$. Trek de translaties $B_{1u} + B_{2u} + B_{3u}$ en de twee rotaties
$B_{2g} + B_{3g}$ af: $2A_g + B_{2g} + B_{3g} + B_{1u} + B_{2u} + B_{3u}$,
zeven trillingen.
In $D_{\infty h}$: $2\Sigma_g^+ + \Pi_g + \Sigma_u^+ + \Pi_u$. IR: $\Sigma_u^+$ en $\Pi_u$
(twee banden); Raman: $2\Sigma_g^+$ en $\Pi_g$ (drie banden); geen
samenvallende banden.
\end{solution}

\begin{solution}{exo:b3:group-theory-applied:12}
$1b_2 \to 4a_1$: $B_2$, $y$, toegestaan, loodrecht op het vlak. $3a_1 \to 4a_1$:
$A_1$, $z$, toegestaan, langs de $C_2$-as. $1b_2 \to 2b_1$: $A_2$, verboden.
$1b_1 \to 4a_1$: $B_1$, $x$, toegestaan, in het vlak, loodrecht op de as.
\end{solution}

\begin{solution}{pb:b3:group-theory-applied:1}
\textbf{1.} cis: de twee CO op naburige hoekpunten; trans: tegenover
elkaar; fac: drie CO op één zijvlak; mer: drie CO op een meridiaan (twee
in trans, één ertussen).
\textbf{2.} $C_{2v}$.
\textbf{3.} $D_{4h}$.
\textbf{4.} $C_{3v}$.
\textbf{5.} $C_{2v}$.
\textbf{6.} Alleen trans-\ce{ML4(CO)2}.
\textbf{7.} Een operatie die een CO op een ander brengt, plaatst zijn
vector buiten de diagonaal (bijdrage 0); een operatie die een CO op zijn
plaats laat, beeldt zijn bindingsvector op zichzelf af (+1); geen enkele
keert een vector C--O om.
\textbf{8.} \omterm{def:b3:point-groups:representation}{Karakters} $(2, 0, 2, 0)$, met $\sigma_v(xz)$ het vlak van beide CO: $A_1 +
B_1$.
\textbf{9.} \omterm{def:b3:point-groups:representation}{Karakters} 2 onder $E$, $C_4$, $C_2$, $\sigma_v$, $\sigma_d$, elders 0:
$A_{1g} + A_{2u}$.
\textbf{10.} $(3, 0, 1)$: $A_1 + E$.
\textbf{11.} $(3, 1, 3, 1)$: $2A_1 + B_1$.
\textbf{12.} 2, 2, 3, 3: één dimensie per CO.
\textbf{13.} \omterm{def:b3:group-theory-applied:activity}{IR-actief}: cis 2, trans 1, fac 2, mer 3.
\textbf{14.} \omterm{def:b3:group-theory-applied:activity}{Raman-actief}: cis 2, trans 1, fac 2, mer 3.
\textbf{15.} trans: zijn IR-band ($A_{2u}$) en zijn ramanband ($A_{1g}$) horen
bij verschillende trillingen.
\textbf{16.} $A_{2u}$ verandert van teken onder $i$: het ene CO wordt
langer terwijl het andere korter wordt, in tegenfase.
\textbf{17.} $E$ is een ontaard paar: twee trillingen met dezelfde
frequentie geven één band, naast de band van $A_1$.
\textbf{18.} Een trans-gedisubstitueerd, centrosymmetrisch isomeer.
\textbf{19.} Drie IR-banden: mer (fac zou er twee geven).
\textbf{20.} Ja: mer heeft drie Raman-actieve CO-trillingen bij dezelfde
frequenties, omdat het niet centrosymmetrisch is; fac zou er twee vertonen.
\textbf{21.} $A_1$: de strekvibratie in fase is totaalsymmetrisch.
\textbf{22.} Een mengsel fac/mer zou tot vijf IR-banden vertonen, met de
banden van fac op andere posities; drie zuivere banden passen bij één
enkel isomeer. Een
${}^{13}$C-NMR-spectrum zou de knoop doorhakken: twee carbonylsignalen in
de verhouding 2:1 voor
mer, één signaal voor fac.
\textbf{23.} \textbf{Het mer-isomeer heeft drie IR-actieve
CO-strekvibraties, het fac-isomeer twee}: het product is mer.
\end{solution}
